\chapter{Livrer : empaqueter, signer, déployer}

Un exécutable dans \texttt{Build/Bin/} n'est pas une livraison. Ce chapitre
couvre les quatre commandes qui produisent quelque chose qu'on donne à quelqu'un
--- et le piège qui a coûté la moitié d'un lot de paquets sur ce dépôt.

\section{\texttt{package}}

\begin{nkcode}{}
jenga package --platform windows --project MonJeu --output dist/
\end{nkcode}

\begin{nkretenir}
\textbf{\texttt{package} CONSTRUIT d'abord.} Il ne se contente pas de zipper ce
qui traîne : il lance la construction, puis empaquette. La fraîcheur est donc
garantie --- ce qu'un empaqueteur pur ne peut pas promettre.

C'est une raison de ne pas écrire son propre script de zippage : une seconde
façon d'empaqueter finit toujours par diverger de la première.
\end{nkretenir}

\begin{nkpiege}[Le piège qui écrase la moitié d'un lot, en silence]
\textbf{\texttt{jenga package} nomme sa sortie d'après le PROJET, pas d'après la
plateforme.}

Windows produit \texttt{MonJeu.zip}. Web produit \texttt{MonJeu.zip}. Dans le
même dossier de sortie, \textbf{le second écrase le premier} --- code de sortie
0, message « package created », et un artefact de moins.

\textbf{Le remède} : un dossier de sortie \emph{par plateforme}, puis un
renommage.

\alerte{} \textbf{Et ce n'est pas la sortie de la commande qui l'a révélé : c'est
l'ouverture de l'archive.} On y a trouvé un \texttt{.wasm} là où l'on attendait
un \texttt{.exe}. Un empaquetage qui « réussit » cinq fois peut ne laisser que
trois artefacts.
\end{nkpiege}

\begin{nkretenir}
\textbf{Deux plateformes ne passent pas par \texttt{package}}, et il faut le
savoir avant d'écrire un script :

\begin{center}
\begin{tabular}{@{}ll@{}}
\toprule
Android & l'APK vient de \texttt{jenga build} \\
HarmonyOS & \texttt{package} exige un \texttt{hvigor-config.json5} absent \\
\bottomrule
\end{tabular}
\end{center}

Un script d'empaquetage a donc deux chemins, et c'est délibéré --- pas une
rustine.
\end{nkretenir}

\section{\texttt{keygen} et \texttt{sign}}

\begin{nkcode}{}
jenga keygen
jenga sign
\end{nkcode}

\begin{nkretenir}
La signature se déclare dans le descripteur, et chaque plateforme a ses champs :

\begin{center}
\begin{tabular}{@{}ll@{}}
\toprule
Windows & \texttt{signingcertificate}, \texttt{signingthumbprint}, \texttt{signingtimestampurl} \\
Android & \texttt{androidkeystore}, \texttt{androidkeyalias}, \texttt{androidkeystorepass} \\
Apple & \texttt{signingidentity}, \texttt{signingentitlements} \\
HarmonyOS & \texttt{harmonycertfile}, \texttt{harmonykeystore}, \texttt{harmonyprofile} \\
\bottomrule
\end{tabular}
\end{center}

\textbf{\texttt{signingtimestampurl} mérite une phrase} : sans horodatage, une
signature devient invalide le jour où le certificat expire --- même pour un
binaire signé bien avant. Avec, elle reste valable.
\end{nkretenir}

\begin{nkpiege}
\textbf{Un APK signé avec la clé de débogage s'installe et ne se publie pas.}
C'est parfait pour des testeurs internes, et refusé par les magasins.

Le script d'empaquetage de ce dépôt le dit à chaque exécution :
\emph{« APK Release signe avec debug.keystore (OK testeurs internes, pas Play
Store) »}.

\emph{Dire ce qu'un artefact n'est pas vaut autant que dire ce qu'il est.}
\end{nkpiege}

\section{\texttt{deploy} et \texttt{publish}}

\begin{center}
\begin{tabular}{@{}ll@{}}
\toprule
\texttt{jenga deploy} & installe sur une cible --- un téléphone, un appareil \\
\texttt{jenga publish} & envoie vers une destination de distribution \\
\bottomrule
\end{tabular}
\end{center}

\begin{nkretenir}[Ces deux-là sortent vers l'extérieur]
\textbf{Toute commande qui envoie quelque chose hors de votre machine mérite une
décision, pas une habitude.} Publier un artefact anodin annonce quand même une
intention, et ce qui est parti est difficile à reprendre.

Ce livre ne les documente pas plus loin parce qu'elles n'ont pas été exercées
ici. \emph{Les décrire d'après leur nom serait exactement ce que le chapitre 12
reproche à un commentaire qui nomme une plateforme.}
\end{nkretenir}

\section{Vérifier un paquet, et non la commande}

\begin{nkretenir}
\textbf{« L'empaquetage a réussi » ne prouve rien.} Le piège de la section 1 le
montre : la commande peut réussir et l'artefact ne pas être celui qu'on croit.

Deux témoins, du plus fort au plus faible :

\textbf{1. Lancer ce qui est DANS le paquet.} Extraire l'exécutable et le faire
tourner --- son banc, s'il en a un. C'est le seul témoin qui porte sur le
comportement.

\textbf{2. Chercher dans le binaire une chaîne qui n'existe que depuis votre
modification.} Un artefact qui la porte a été construit après. Une \emph{date},
elle, se préserve à la copie et ne prouve rien.

\alerte{} \textbf{Le second témoin exige un contrôle positif} : cherchez d'abord une
chaîne dont vous savez qu'elle est là. Si vous ne la trouvez pas, votre lecteur
ne sait pas lire ce format --- et l'absence de la première ne prouve rien du
tout.
\end{nkretenir}

\begin{nkretenir}
\textbf{Cette méthode a été appliquée sur ce dépôt}, sur quinze paquets et cinq
plateformes. Résultat : les trois exécutables Windows passent leur banc, et les
quatre autres plateformes portent la chaîne témoin.

Et le vérificateur lui-même s'est trompé au premier essai : il comptait le
\emph{mot} « ECHEC » et accusait un jeu de deux échecs --- qui étaient deux
\emph{noms de cas} contenant ce mot. \emph{Compter un mot n'est pas mesurer une
chose}, y compris quand c'est vous qui écrivez l'instrument.
\end{nkretenir}

\section{Ce qui accompagne un paquet}

\begin{nkretenir}
Un dossier de livraison porte plus que des artefacts : il porte de quoi s'en
servir. Le lot de ce dépôt contient un \texttt{LISEZMOI.md} avec un tableau
« fichier / pour qui / comment on s'en sert », et un avertissement en tête :

% `escapeinside` : la regle `literate` du preambule ne mord PAS ici -- mesure
% faite, y compris apres avoir mis `language={}`, donc ce n'est pas le lexer.
% `escapeinside` ne depend d'aucune analyse : ce qui est entre `(*` et `*)` est
% du LaTeX ordinaire. C'est le seul mecanisme qui ait fonctionne.
\begin{nkcode}[listing options={escapeinside={(*}{*)}}]{Build/JeuxPlateau-Release/LISEZMOI.md}
## (*\alerte{}*) Le Web ne s'ouvre PAS en double-cliquant le `.html`
\end{nkcode}

\textbf{Cette phrase économise plus de messages que toute la documentation
technique du lot.} Écrivez celle qui correspond à \emph{votre} piège le plus
probable.
\end{nkretenir}

\begin{nkbilan}
\texttt{package} construit avant d'empaqueter, ce qui garantit la fraîcheur ---
mais il nomme sa sortie d'après le projet, et deux plateformes s'écrasent en
silence dans un même dossier. Android et HarmonyOS passent par \texttt{build},
pas par \texttt{package}. La signature se déclare par plateforme, et une clé de
débogage donne un paquet installable mais non publiable --- il faut le dire.
Enfin, un paquet se vérifie en \emph{lançant ce qu'il contient}, pas en lisant la
sortie de la commande, et tout témoin par recherche de chaîne exige son contrôle
positif.
\end{nkbilan}

\begin{nkquestions}
\begin{enumerate}
\item Que fait \texttt{jenga package} avant d'empaqueter ?
\item D'après quoi nomme-t-il sa sortie, et quelle conséquence ?
\item Quelles plateformes ne passent pas par \texttt{package} ?
\item À quoi sert \texttt{signingtimestampurl} ?
\item Citez deux témoins de la fraîcheur d'un paquet, et lequel est le plus fort.
\end{enumerate}
\end{nkquestions}

\begin{nkpratique}
Empaquetez votre projet pour deux plateformes \textbf{dans le même dossier}, à
dessein.

\begin{enumerate}
\item comptez les fichiers produits ;
\item ouvrez celui qui reste et regardez ce qu'il contient ;
\item recommencez avec un dossier par plateforme, puis un renommage ;
\item comparez les deux comptes.
\end{enumerate}

Le point 2 est le c\oe ur de l'exercice : c'est en ouvrant qu'on voit, jamais
en lisant la sortie de la commande.
\end{nkpratique}

\begin{nkdemo}
Prouvez qu'un paquet contient bien vos dernières modifications.

\begin{enumerate}
\item ajoutez à votre programme une chaîne unique --- la date du jour ;
\item construisez et empaquetez ;
\item extrayez l'exécutable du paquet et \textbf{lancez-le} : il doit afficher
      la chaîne ;
\item cherchez ensuite cette même chaîne dans le binaire, sans le lancer.
\end{enumerate}

\textbf{Ce que la démonstration établit} : les deux témoins concordent. Faites
alors le contraire --- empaquetez \emph{sans} reconstruire après une nouvelle
modification --- et vérifiez qu'ils sont tous deux capables de dire non.

\alerte{} Un témoin qui n'a jamais dit non n'a rien prouvé.
\end{nkdemo}

\begin{nklogique}
Vous livrez cinq paquets. Un testeur rapporte que « le jeu n'a pas la correction
d'hier ».

\begin{enumerate}
\item Énumérez quatre causes, depuis « la correction n'est pas dans le dépôt »
      jusqu'à « le testeur a installé l'ancien ».
\item Pour chacune, dites qui peut la vérifier --- vous, ou lui --- et avec quoi.
\item Une seule de ces causes se vérifie \emph{sans} le testeur et \emph{sans}
      reconstruire. Laquelle, et pourquoi commencer par elle ?
\end{enumerate}
\end{nklogique}
